
\subsection{Materialabhängigkeit des Widerstandes R\label{Materialabhängigkeit des Widerstandes}}
\Aufgabe{
    Ein Draht aus Kupfer hat folgende Eigenschaften:
    \begin{itemize}
        \item Länge: $ \ell = 5\,\mathrm{m} $
        \item Querschnittsfläche: $ A = 1\,\mathrm{mm}^2$
        \item spezifischer Widerstand von Kupfer: $\rho = 1,68\cdot 10^{-8}\,\Omega\mathrm{m}$
    \end{itemize}
    Berechne $ R $ des Drahtes.
}
\Loesung{
    \begin{equation*}
        R = \rho \cdot \frac{\ell}{A} = 1,68 \cdot 10^{-8} \cdot \frac{5}{1 \cdot 10^{-6}} \,\Omega= 0,084\,\Omega
    \end{equation*}
}



\subsection{Elektrischer Widerstand \label{Elektrischer Widerstand}}
\Aufgabe{
    Eine kugelförmige Hohlstruktur besteht aus zwei ideal leitfähigen Schalen – einer inneren mit Radius $ a $ und einer äußeren mit Radius $ b $.
    Der Zwischenraum ist vollständig mit einem homogenen, leitfähigen Material gefüllt, das die elektrische Leitfähigkeit $ \kappa $ (in $\frac{\mathrm{S}}{\mathrm{m}}$) besitzt.
    Zwischen den leitenden Schalen wird durch eine Spannungsquelle eine Spannung
    \begin{equation*}
        U = \varphi_1 - \varphi_2
    \end{equation*}
    angelegt.
    Der Strom kann sich aufgrund der vollständigen Füllung und der idealen Kontakte ungehindert ausbreiten.

    Berechnen Sie den elektrischen Widerstand $ R $ zwischen den beiden kugelförmigen Schalen in Abhängigkeit von $ a $, $ b $ und $ \kappa $.
}
\Loesung{
    Da das System radial-symmetrisch ist, fließt der Strom radial von innen nach außen (oder umgekehrt), und das elektrische Feld hängt nur vom Radius $ r $ ab.
    Die Stromdichte ergibt sich über das Ohmsche Gesetz in Differentialform:
    \begin{equation*}
        \vec{j}(r) = \kappa \cdot \vec{E}(r)
    \end{equation*}

    Die Fläche, durch die der Strom bei Radius $ r $ fließt, ist die Oberfläche einer Kugel:
    \begin{equation*}
        A(r) = 4 \pi r^2
    \end{equation*}

    Der Gesamtstrom ist konstant durch alle Kugelflächen:
    \begin{equation*}
        I = j(r) \cdot A(r) = \kappa \cdot E(r) \cdot 4 \pi r^2 \quad \Rightarrow \quad E(r) = \frac{I}{4 \pi \kappa r^2}
    \end{equation*}

    \textbf{Spannung berechnen:}

    Die Spannung ergibt sich durch das Integral über das elektrische Feld:
    \begin{equation*}
        U = \int_a^b E(r) \, dr = \int_a^b \frac{I}{4 \pi \kappa r^2} \, dr = \frac{I}{4 \pi \kappa} \int_a^b \frac{1}{r^2} \, dr = \frac{I}{4 \pi \kappa} \left[-\frac{1}{r}\right]_a^b
    \end{equation*}
    \begin{equation*}
        U = \frac{I}{4 \pi \kappa} \left(\frac{1}{a} - \frac{1}{b}\right)
    \end{equation*}

    \textbf{Widerstand berechnen:}

    Da $ U = R \cdot I $ gilt, folgt
    \begin{equation*}
        R = \frac{U}{I} = \frac{1}{4 \pi \kappa} \left(\frac{1}{a} - \frac{1}{b}\right)
    \end{equation*}

    Der elektrische Widerstand zwischen den beiden Schalen beträgt somit:
    \begin{equation*}
        \boxed{ R = \frac{1}{4 \pi \kappa} \left(\frac{1}{a} - \frac{1}{b}\right) }
    \end{equation*}
}




\subsection{Elektrischer Widerstand eines beschichteten Drahts\label{Elektrischer Widerstand eines beschichteten Drahts}}
\Aufgabe{
    Ein runder Aluminiumdraht besitzt einen Radius von $ a = 0,5\,\mathrm{mm}$. Er ist mit einer dünnen Goldschicht von $ b = 0,1\,\mathrm{mm}$ Dicke überzogen.
    Für die spezifische elektrische Leitfähigkeit und den Temperaturkoeffizienten der beiden Metalle gelten folgende Werte:

    \begin{itemize}
        \item \textbf{Aluminium:}
            \begin{equation*}
                \kappa_\text{Alu} = 35\,\frac{\mathrm{m}}{\Omega\mathrm{mm}^2}, \quad \alpha_\text{Alu} = 3,5 \cdot 10^{-3}\,\frac{1}{\mathrm{K}}
            \end{equation*}
        \item \textbf{Gold:}
            \begin{equation*}
                \kappa_\text{Gold} = 44\,\frac{\mathrm{m}}{\Omega\mathrm{mm}^2}, \quad \alpha_\text{Gold} = 3{,}6 \cdot 10^{-3}\,\frac{1}{\mathrm{K}}
            \end{equation*}
    \end{itemize}

    \textbf{Gegeben:}
    Drahtlänge: $ l = 1,5\,\mathrm{m} $
    Temperatur: $ T = 25\,\degree\mathrm{C} $ (Ausgangstemperatur)

    \begin{itemize}
      \item[a)] Berechnen Sie den elektrischen Gleichstromwiderstand des Aluminiumkerns und der Goldummantelung bei $ T = 25\,\degree\mathrm{C} $.
      \item[b)] Wie verändert sich der jeweilige Gleichstromwiderstand, wenn die Temperatur auf $ T = 90\,\degree\mathrm{C} $ ansteigt?
    \end{itemize}
}
\Loesung{
    \begin{align*}
        \text{Radius Aluminiumkern:} &\quad a = 0,5\,\mathrm{mm} \\
        \text{Dicke der Goldschicht:} &\quad b = 0,1\,\mathrm{mm} \\
        \Rightarrow \text{Gesamtradius:} &\quad r = a + b = 0,6\,\mathrm{mm} \\
        \text{Länge des Leiters:} &\quad l = 1,5\mathrm{m} = 1500\,\mathrm{mm} \\
        \text{Aluminium:} \quad \kappa_\text{Alu} = 35\,\frac{\mathrm{m}}{\Omega\mathrm{mm}^2}, &\quad \alpha_\text{Alu} = 3,5 \cdot 10^{-3}\, \frac{1}{\mathrm{K}} \\
        \text{Gold:} \quad \kappa_\text{Gold} = 44\,\frac{\mathrm{m}}{\Omega\mathrm{mm}^2}, &\quad \alpha_\text{Gold} = 3{,}6 \cdot 10^{-3}\, \frac{1}{\mathrm{K}} \\
        \text{Temperaturänderung:} &\\
        &T_1 = 25\,\degree\mathrm{C} = 298\,\mathrm{K} \\
        &T_2 = 90\,\degree\mathrm{C} = 363\,\mathrm{K} \\
        &\Delta T = T_2 - T_1 = 65\,\mathrm{K}
    \end{align*}

    Querschnittsflächen berechnen

    \begin{align*}
        &\text{Aluminiumkern:} \quad A_\text{Alu} = \pi a^2 = \pi \cdot 0,5^2 = \pi \cdot 0,25 = 0,7854\,\mathrm{mm}^2 \\
        &\text{Goldmantel (Ringfläche zwischen } r = 0,6\,\mathrm{mm} \text{ und } a = 0,5\,\mathrm{mm):} \\
        &A_\text{Gold} = \pi (r^2 - a^2) = \pi ((0,6\,\mathrm{mm})^2 - (0,5\,\mathrm{mm})^2) = \pi (0,36\,\mathrm{mm}^2 - 0,25\,\mathrm{mm}^2) = \pi \cdot 0,11\,\mathrm{mm}^2 \approx 0,3456\,\mathrm{mm}^2
    \end{align*}

    Widerstände bei $25\,\degree\mathrm{C}$

    Formel:
    \begin{equation*}
        R = \frac{l}{\kappa \cdot A}
    \end{equation*}

    \begin{align*}
        R_\text{Alu} &= \frac{1500}{35 \cdot 0,7854} \,\Omega= \frac{1500}{27,489} \,\Omega\approx 54,58\,\mathrm{m}\Omega \\
        R_\text{Gold} &= \frac{1500}{44 \cdot 0,3456} \,\Omega= \frac{1500}{15,206} \,\Omega\approx 98,61\,\mathrm{m}\Omega
    \end{align*}

    Widerstände bei $90\,\degree\mathrm{C}$

    Temperaturabhängiger Widerstand:
    \begin{equation*}
        R(T) = R_0 \cdot \left( 1 + \alpha \cdot \Delta T \right)
    \end{equation*}

    \begin{align*}
        R_{\text{Alu},90^\circ \mathrm{C}} &= 54,58 \cdot (1 + 3,5 \cdot 10^{-3} \cdot 65) \,\Omega= 54,58 \cdot 1,2275 \,\Omega\approx 66,97\,\mathrm{m}\Omega \\
        R_{\text{Gold},90^\circ \mathrm{C}} &= 98,61 \cdot (1 + 3,6 \cdot 10^{-3} \cdot 65)\,\Omega = 98,61 \cdot 1,234 \,\Omega\approx 121,64\,\mathrm{m}\Omega
    \end{align*}


    \begin{tabular}{lcc}
        \toprule
            Material & Widerstand bei $25\,\degree\mathrm{C}$ & Widerstand bei $90\,\degree\mathrm{C}$ \\
        \midrule
            Aluminium & $54,58\,\mathrm{m}\Omega$ & $66,97\,\mathrm{m}\Omega$ \\
            Gold & $98,61\,\mathrm{m}\Omega$ & $121,64\,\mathrm{m}\Omega$ \\
        \bottomrule
    \end{tabular}
}





\subsection{Reale Stromquellen}
\Aufgabe{
    Eine reale Stromquelle hat eine Leerlaufspannung $ U_0 = 9\,\mathrm{V} $ und einen Innenwiderstand $ R_\text{i} = 1\,\Omega $.

    a) wie groß ist die Klemmenspanung, wenn ein Lastwiderstand von  $ R_l = 4\,\Omega $ angeschlossen wird ?
    b) wie groß ist der Strom durch die Last?
}
\Loesung{
    \textbf{a) Klemmenspannung $ U_\text{k} $:}
    Die Klemmenspannung berechnet sich mit dem Spannungsteiler:
    \begin{equation*}
        U_\text{k} = U_0 \cdot \frac{R_\text{l}}{R_\text{i} + R_\text{l}} = 9\,\mathrm{V} \cdot \frac{4\,\Omega}{1\,\Omega + 4\,\Omega} = 9 \cdot \frac{4}{5} = 7,2\,\mathrm{V}
    \end{equation*}

    \textbf{b) Strom durch die Last $ I $:}

    \begin{equation*}
        I = \frac{U_0}{R_\text{i} + R_\text{l}} = \frac{9\,\mathrm{V}}{1\,\Omega + 4\,\Omega} = \frac{9\,\mathrm{V}}{5\,\Omega} = 1,8\,\mathrm{A}
    \end{equation*}
}





\subsection{Spannungsquelle}
\Aufgabe{
    Gegeben ist ein Netzwerk aus einer Spannungsquelle mit $ U = 10\,\mathrm{V} $,
    einem Reihenwiderstand $ R_1 = 2\,\Omega $
    und einem parallelen Widerstand $ R_2 = 4\,\Omega $.

    Berechne Ströme und Spannungen im Netzwerk.
}
\Loesung{
    \begin{equation*}
        R_\text{ges} = R_1 + R_2 = 2\,\Omega + 4\,\Omega = 6\,\Omega, \quad
        I = \frac{U}{R_\text{ges}} = \frac{10\,\mathrm{V}}{6\,\Omega} = 1,67\,\mathrm{A}
    \end{equation*}

    \begin{equation*}
        U_{R_1} = I \cdot R_1 = 1,67\,\mathrm{A} \cdot 2\,\Omega = 3,33\,\mathrm{V}, \quad
        U_{R_2} = U - U_{R_1} = 10\,\mathrm{V} - 3,33\,\mathrm{V} = 6,67\,\mathrm{V}
    \end{equation*}

    \begin{equation*}
        I_{R_2} = \frac{U_{R_2}}{R_2} = \frac{6,67\,\mathrm{V}}{4\,\Omega} = 1,67\,\mathrm{A}
    \end{equation*}
}






\subsection{Kapazität und Kondensator\label{Kapazität und Kondensator}}
\Aufgabe{
\textbf{Teil 1: Einzelner Kondensator}
    Gegeben:
    \begin{equation*}
        C = 10\,\mathrm{pF}, \quad U = 5\,\mathrm{V}
    \end{equation*}

    \begin{itemize}
        \item[a)] Berechne die Ladung $ Q $
        \item[b)] Berechne die im Kondensator gespeicherte Energie
    \end{itemize}

    \vspace{1em}

    \textbf{Teil 2: Zwei Kondensatoren in Reihe}

    Zwei Kondensatoren $ C_1 $ und $ C_2 $ sind in Reihe geschaltet.

    \begin{equation*}
        C_1 = 4\,\mathrm{\mu F}, \quad C_2 = 6\,\mathrm{\mu F}, \quad U_\text{ges} = 12\,\mathrm{V}
    \end{equation*}

    \begin{itemize}
        \item[c)] Berechne die Gesamtkapazität $ C_\text{ges} $
        \item[d)] Wie groß ist die Spannung über $ C_1 $?
    \end{itemize}
}
\Loesung{
    \begin{enumerate}
        \item[a)] \textbf{Ladung:}
            \begin{equation*}
                Q = C \cdot U = 10 \cdot 10^{-12} \cdot 5 \,\mathrm{C}= 50\,\mathrm{pC}
            \end{equation*}

        \item[b)] \textbf{Gespeicherte Energie:}
            \begin{equation*}
                E = \frac{1}{2} C U^2 = \frac{1}{2} \cdot 10 \cdot 10^{-12} \cdot 25\,\mathrm{J} = 125\,\mathrm{pJ}
            \end{equation*}

        \item[c)] \textbf{Gesamtkapazität der Reihenschaltung:}
            \begin{equation*}
                \frac{1}{C_\text{ges}} = \frac{1}{C_1} + \frac{1}{C_2} = \frac{1}{4\,\mathrm{\mu F}} + \frac{1}{6\,\mathrm{\mu F}} = \frac{5}{12\,\mathrm{\mu F}}
                \Rightarrow C_\text{ges} = 2,4\,\mathrm{\mu F}
            \end{equation*}

        \item[d)] \textbf{Spannung über $ C_1 $:}

        \begin{equation*}
            Q = C_\text{ges} \cdot U_\text{ges} = 2,4 \cdot 10^{-6} \cdot 12 \,\mathrm{C}= 28,8\,\mathrm{\mu C}
        \end{equation*}
        \begin{equation*}
            U_1 = \frac{Q}{C_1} = \frac{28,8 \cdot 10^{-6}}{4 \cdot 10^{-6}} \,\mathrm{V}= 7,2\,\mathrm{V}
        \end{equation*}
    \end{enumerate}
}





\subsection{Induktivität und Spule\label{Induktivität und Spule}}
\Aufgabe{
    \textbf{Teil 1: Energie in der Spule}
    Gegeben ist eine Spule mit der Induktivität
    \begin{equation*}
        L = 2\,\mathrm{H}
    \end{equation*}
    und einem Strom
    \begin{equation*}
        I = 3\,\mathrm{A}.
    \end{equation*}

    \begin{itemize}
        \item[a)] Berechne die im Magnetfeld der Spule gespeicherte Energie.
        \item[b)] Wie ändert sich die Energie, wenn der Strom auf $ I = 0,5\,\mathrm{A} $ reduziert wird?
    \end{itemize}

    \textbf{Teil 2: Induktivität einer zylindrischen Spule}

    Eine zylindrische Spule hat
    \begin{equation*}
        N = 500, \quad l = 20\,\mathrm{cm}, \quad A = 5\,\mathrm{cm}^2
    \end{equation*}
    Das Innere der Spule ist luftgefüllt.

    Berechne die Induktivität $ L $ der Spule.
}
\Loesung{
    \begin{enumerate}
    \item[a)] \textbf{Energie bei $ I = 3\,\mathrm{A} $:}
        \begin{equation*}
            E = \frac{1}{2} L I^2 = \frac{1}{2} \cdot 2 \cdot 3^2 \,\mathrm{J}= 9\,\mathrm{J}
        \end{equation*}

    \item[b)] \textbf{Energie bei $ I = 0,5\,\mathrm{A} $:}

    \begin{equation*}
        E = \frac{1}{2} \cdot 2 \cdot 0,5^2\,\mathrm{J} = 0,25\,\mathrm{J}
    \end{equation*}

    \item[c)] \textbf{Berechnung der Induktivität $ L $ der zylindrischen Spule:}

    Zunächst die Umrechnung der Einheiten:
    \begin{equation*}
        \ell = 20\,\mathrm{cm} = 0,2\,\mathrm{m}, \quad A = 5\,\mathrm{cm}^2 = 5 \cdot 10^{-4}\,\mathrm{m}^2
    \end{equation*}

    Die Induktivität einer luftgefüllten zylindrischen Spule berechnet sich durch:
    \begin{equation*}
        L = \mu_0 \frac{N^2 A}{\ell}
    \end{equation*}
    mit der magnetischen Feldkonstanten
    \begin{equation*}
        \mu_0 = 4\pi \cdot 10^{-7}\,\frac{\mathrm{H}}{\mathrm{m}}
    \end{equation*}

    Einsetzen:
    \begin{equation*}
        L = 4\pi \cdot 10^{-7} \cdot \frac{500^2 \cdot 5 \cdot 10^{-4}}{0,2} \,\mathrm{H}= 7,85 \cdot 10^{-2} \,\mathrm{H} = 78,5\,\mathrm{mH}
    \end{equation*}
    \end{enumerate}
}










\subsection{Induktivität, Magnetische Flussdichte und Magnetfeldenergie\label{Induktivität, Magnetische Flussdichte und Magnetfeldenergie}}
\Aufgabe{
    Ein luftgefüllter, langer zylindrischer Spulenleiter besitzt folgende Eigenschaften:

    \begin{itemize}
      \item Anzahl der Windungen: $ N = 500 $
      \item Länge der Spule: $ l = 0,5\,\mathrm{m} $
      \item Querschnittsfläche: $ A = 4 \cdot 10^{-4}\,\mathrm{m}^2 $
      \item Stromstärke: $ I = 2,5\,\mathrm{A} $
    \end{itemize}

    \textbf{Gegeben:}
    Magnetische Feldkonstante: $ \mu_0 = 4\pi \cdot 10^{-7}\,\frac{\mathrm{H}}{\mathrm{m}} $

    \begin{itemize}
      \item[a)] Berechnen Sie die magnetische Flussdichte $ \vec{B} $ im inneren der Spule.
      \item[b)] Bestimmen Sie die Induktivität $ L $ der Spule.
      \item[c)] Berechnen Sie die im Magnetfeld gespeicherte Energie $ E_m $.
    \end{itemize}

}
\Loesung{
    \textbf{a) Magnetische Flussdichte $ \vec{B} $:}

    Für eine lange, luftgefüllte Zylinderspule gilt:
    \begin{equation*}
        B = \mu_0 \cdot \frac{N}{l} \cdot I
    \end{equation*}

    Einsetzen der Werte:
    \begin{equation*}
        B = 4\pi \cdot 10^{-7} \cdot \frac{500}{0,5} \cdot 2,5 \,\mathrm{T}= 4\pi \cdot 10^{-7} \cdot 1000 \cdot 2,5\,\mathrm{T}
    \end{equation*}
    \begin{equation*}
        B = 4\pi \cdot 10^{-7} \cdot 2500 \,\mathrm{T}= \pi \cdot 10^{-3}\,\mathrm{T}
    \end{equation*}
    \begin{equation*}
        \Rightarrow B \approx 3,14 \cdot 10^{-3}\,\mathrm{T} = 3,14\,\mathrm{mT}
    \end{equation*}

    \textbf{b) Induktivität $ L $:}

    Die Induktivität einer langen luftgefüllten Spule ergibt sich aus:
    \begin{equation*}
        L = \mu_0 \cdot \frac{N^2 \cdot A}{l}
    \end{equation*}

    Einsetzen der Werte:
    \begin{equation*}
        L = 4\pi \cdot 10^{-7} \cdot \frac{500^2 \cdot 4 \cdot 10^{-4}}{0,5} \,\mathrm{H}
    \end{equation*}
    \begin{equation*}
        L = 4\pi \cdot 10^{-7} \cdot \frac{250000 \cdot 4 \cdot 10^{-4}}{0,5}\,\mathrm{H}
        = 4\pi \cdot 10^{-7} \cdot \frac{100}{0,5}\,\mathrm{H}
        = 4\pi \cdot 10^{-7} \cdot 200 \,\mathrm{H}
    \end{equation*}
    \begin{equation*}
        L = 8\pi \cdot 10^{-5} \,\mathrm{H}
        \Rightarrow L \approx 2,51\cdot 10^{-4}\,\mathrm{H} = 251\,\mathrm{\mu H}
    \end{equation*}

    \textbf{c) Gespeicherte Energie im Magnetfeld $ E_m $:}

    Die Energie im Magnetfeld einer Spule berechnet sich mit:
    \begin{equation*}
        E_m = \frac{1}{2} L I^2
    \end{equation*}

    Einsetzen:
    \begin{equation*}
        E_m = \frac{1}{2} \cdot 2,51 \cdot 10^{-4} \cdot (2,5)^2\,\mathrm{J}
        = 0,5 \cdot 2,51 \cdot 10^{-4} \cdot 6,25\,\mathrm{J}
    \end{equation*}
    \begin{equation*}
        E_m = 7,84 \cdot 10^{-4}\,\mathrm{J}
        = 0,784\,\mathrm{mJ}
    \end{equation*}
}
